\subsection{$C^k$-convergence}

Let $X,Y$ be smooth manifolds. In this section we discuss the concept of $C^k$-convergence for sequences of functions $(f_n) \subseteq C^k(X,Y)$ with respect to the Whitney topolgies. We will be using the definition of the Whitney topolgies based on jet bundles. The notation and terminology used will be the same as in \cite{golubitsky2012stable}. First we remark that $X$ or $Y$ is also allowed to have boundary. The definition of the jet bundles $J^k(X,Y)$ and the Whitney topologies works exactly the same \footnotemark. Let us begin by understanding the $C^0$-topology.

\footnotetext{If both have boundary, $J^k(X,Y)$ will be a manifold with corners, so we might want to avoid that. However $J^k(X,Y)$ still carries a well-defined topology and I think this is all we need and the proofs in this section all go through the same. In \cite{hirsch2012differential} at the end of Chapter 2.4 on page 63 the author briefly describes the definition of $J^k(X,Y)$ in this case.}

\begin{lemma}
    \label{whitney_top:lem:c0_and_exponential_topology_agree}
    Let $X,Y$ be smooth manifolds with $X$ compact. Then the Whitney $C^0$-topology (as described in \cite{golubitsky2012stable}) and the exponential topology (as described in \cite{steimle_topo}) on $C^0(X,Y)$ agree.
\end{lemma}

\begin{proof}
    We identify $J^0(X,Y) \cong X \times Y$. Then by definition the sets of the form 
    \[
    \mathrm{M}(W) := \{f \in C^0(X,Y) \mid j^0f(X) = \mathrm{Graph}(f) \subseteq W\}
    \]
    for $W\subseteq X \times Y$ open are a basis of the $C^0$-topology on $C^0(X,Y)$. To see that this topology agrees with the exponential topology we show that it satisfies the correct universal property. So let $S$ be any topological space. We have to show that a map 
    \[
    F \colon S \to C^0(X,Y)
    \]
    is continuous if and only if 
    \[
    \widetilde{F} \colon S \times X \to Y, \; (s,x) \mapsto F(s)(x)
    \]
    is continuous. First assume $F$ to be continuous. Let $V\subseteq Y$ be open and $(s_0,x_0) \in \widetilde{F}^{-1}(V)$. Define $f_0 := F(s_0) \in C^0(X,Y)$, then 
    \[
    U := f_0^{-1}(V)
    \]
    is open in $X$. and $x_0 \in U$. Choose some smaller open neighbourhood $U_0 \subseteq X$ of $x_0$ with $\overline{U}_0 \subseteq U$ and set 
    \[
    W := U \times V \cup (X\setminus \overline{U}_0) \times Y \subseteq X \times Y.
    \]
    By assumption $S_0 := F^{-1}(\mathrm{M}(W)) \subseteq S$ is open and we have constructed $W$, such that $s_0 \in S_0$. But then 
    \[
    S_0 \times U_0 \subseteq S \times X
    \]
    is an open neighbourhood of $(s_0,x_0)$ and $S_0 \times U_0 \subseteq \widetilde{F}^{-1}(V)$.

    Now assume that $\widetilde{F}$ is continuous. Let $W\subseteq X \times Y$ be open and $s_0 \in F^{-1}(\mathrm{M}(W))$. Again define $f_0 := F(s_0) \in C^0(X,Y)$. Then $\mathrm{Graph}(f_0) \subseteq W$ and around each $(x,f_0(x)) \in \mathrm{Graph}(f_0)$ we can find a product neighbourhood $U_x \times V_x \subseteq W$, where $U_x \subseteq X$ is an open neighbourhood of $x$ and $V_x \subset Y$ is an open neighbourhood of $f_0(x)$. After possibly shrinking $U_x$ we may further assume that $\overline{U}_x \subseteq f_0^{-1}(V_x)$. Since $X$ is compact we can cover $\mathrm{Graph}(f_0)$ with finitely many such neighbourhoods $U_1 \times V_1,\ldots,U_n\times V_n$. Now for each $V_i$ we have that 
    \[
    \{s_0\} \times \overline{U}_i \subseteq \widetilde{F}^{-1}(V_i).
    \]
    But $\widetilde{F}^{-1}(V_i) \subseteq S \times X$ is open by assumption and $\overline{U}_i$ is compact, so we can find open neighbourhoods $S_i \subseteq S$ of $s_0$ with 
    \[
    S_i \times \overline{U}_i \subseteq \widetilde{F}^{-1}(V_i).
    \]
    Set $S_0 := S_1 \cap \ldots \cap S_n$. Then $S_0 \subseteq S$ is an open neighbourhood of $s_0$ and for any $(s,x) \in S_0 \times X$ there exists $U_i$ with $x \in U_i$, hence 
    \[
    F(s)(x) = \widetilde{F}(s,x) \in V_i.
    \]
    It follows that $\mathrm{Graph}(F(s)) \subseteq \bigcup_{i=1}^n U_i \times V_i \subseteq W$ for all $s \in S_0$ and thus $S_0 \subseteq F^{-1}(\mathrm{M}(W))$.
\end{proof}

\begin{remark}
    If $X$ is not compact the $C^0$-topology will in general be finer than the exponential topology. To see this recall that a subbase for the exponential topology on $C^0(X,Y)$ is given by sets of the form
    \[
    F(K,V) := \{f \in C^0(X,Y) \mid f(K) \subseteq V\}
    \]
    for compact sets $K \subseteq X$ and open sets $V \subseteq Y$. Set 
    \[
    W := X \times V \cup (X\setminus K) \times Y.
    \]
    Then $W \subseteq X \times Y$ is open and 
    \[
    F(K,V) = M(W).
    \]
    Before giving an example where the $C^0$-topology is strictly finer, we fix some terminology to differentiate between the two topologies on $C^0(X,Y)$. We shall use the term $C^0_\mathrm{exp}$-convergence to refer to convergence with respect to the exponential topology and reserve the term $C^0$-convergence for convergence with respect to the Whitney $C^0$-topology. Now consider the following example. Define a sequence 
    \[
    f_n := \mathrm{const}_{\frac{1}{n}}  \colon \R \to \R, \; x \mapsto \frac{1}{n}
    \]
    of constant continuous functions. We claim that 
    \[
    f_n \xrightarrow{C^0_\mathrm{exp}} \mathrm{const}_0.
    \]
    To see this let $F(K,V) \subseteq C^0(\R,\R)$ be any open neighbourhood of $\mathrm{const}_0$ of the form described above. Then $0 \in V$ and so for $n$ sufficiently large $\frac{1}{n} \in V$, which is equivalent to $f_n \in F(K,V)$. Since the sets of the form $F(K,V)$ are a subbase for the exponential topology, this shows the claim. However the sequence $(f_n)$ does not $C^0$ converge to $\mathrm{const}_0$, since this would necessarily require for $f_n$ to equal $0$ outside a compact subset for sufficiently large $n$ (see \cite{golubitsky2012stable}).
\end{remark}

Next we understand convergence in the $C^k$-topology. Before stating the main proposition we shall introduce another type of convergence which is better suited for non-compact domains, since in the non-compact case usual $C^k$-convergence is often too strong, because (as already statetd above) 
\[
f_n \xrightarrow{C^k} f
\]
would already imply that $f_n$ and $f$ agree outside a compact subset (again see \cite{golubitsky2012stable}).

\begin{definition}
    Let $X,Y$ be smooth manifolds and $f,f_n \in C^k(X,Y)$. We say that 
    \[
    f_n \xrightarrow{C^k_{\mathrm{loc}}} f
    \]
    if for every compact codimension $0$-submanifold $K \subseteq X$ we have 
    \[
    f_n|_K \xrightarrow{C^k} f|_K.
    \]
\end{definition}

Note that for compact domains $C^k_{\mathrm{loc}}$-convergence and $C^k$-convergence coincide, so we are actually interested in understanding $C^k_{\mathrm{loc}}$-convergence. The following two propositions give a useful characterization of $C^k_{\mathrm{loc}}$-convergence to work with in practice. 

\begin{proposition}
    \label{whitney_top:prop:ck_conv_on_rn}
    Let $D\subseteq \R^n$ be a compact codimension $0$ submanifold, $V \subseteq \R^m$ open and $f,f_n \in C^k(D,V)$. Then
    \[
    f_n \xrightarrow{C^k} f
    \]
    if and only if for every multiindex $\alpha \in \N^n_0$ with $0\leq\abs{\alpha} \leq k$
    \[
    \frac{\partial^{\abs{\alpha}}f_n}{\partial x^\alpha} \xrightarrow{\mathrm{uniformly}}\frac{\partial^{\abs{\alpha}}f}{\partial x^\alpha}.
    \]
\end{proposition}

\begin{proposition}
    \label{whitney_top:prop:ck_conv_equivalent_characterizations}
    Let $X,Y$ be smooth manifolds and $f,f_n \in C^k(X,Y)$. Then the following statements are equivalent.
    \begin{enumerate}[(i)]
        \item $f_n \xrightarrow{C^k_{\mathrm{loc}}} f$,
        \item $j^kf_n \xrightarrow{C^0_{\mathrm{loc}}} j^kf$,
        \item For all charts $\phi \colon U \to U^\prime \subseteq X$ and $\psi \colon V \to V^\prime \subseteq Y$ with $f(U^\prime), f_n(U^\prime) \subseteq V^\prime$ (for $n$ sufficiently large) we have 
        \[
        \psi^{-1} \circ f_n \circ \phi \xrightarrow{C^k_{\mathrm{loc}}} \psi^{-1} \circ f \circ \phi,
        \]
        \item For every $p \in X$ there exist charts $\phi \colon U \to U^\prime \subseteq X$ around $p$ and $\psi \colon V \to V^\prime \subseteq Y$ around $f(p)$ with $f(U^\prime), f_n(U^\prime) \subseteq V^\prime$ (for $n$ sufficiently large) such that
        \[
        \psi^{-1} \circ f_n \circ \phi \xrightarrow{C^k_{\mathrm{loc}}} \psi^{-1} \circ f \circ \phi.
        \]
    \end{enumerate}
\end{proposition}

Before proving the propositions we recall a few facts about the exponential topology not mentioned in \cite{steimle_topo}, which will be used in the proofs.

\begin{lemma}
    \label{whitney_top:lem:exponential_topology_facts}
    Let $X,Y,Z$ be topological spaces and $A \subseteq X$, $B \subseteq Y$ subspaces. Then the following holds (assuming the exponential topology exists on each space)
    \begin{enumerate}[(i)]
        \item The canonical map $C^0(X,B) \hookrightarrow C^0(X,Y)$ is an embedding
        \item The restriction map $C^0(X,Y) \to C^0(A,Y)$ is continuous
        \item The canonical bijection $C^0(X,Y) \times C^0(X,Z) \to C^0(X,Y \times Z)$ is a homeomorphism
    \end{enumerate}
\end{lemma}

\begin{proof}
    For (i) we show that the exponential topology on $C^0(X,B)$ agrees with the topology induced by the map  $C^0(X,B) \hookrightarrow C^0(X,Y)$ by showing the corresponding universal property. So let
    \[\begin{tikzcd}
	& {C^0(X,Y)} \\
	Z & {C^0(X,B)}
	\arrow["{F}", from=2-1, to=1-2]
	\arrow["f", from=2-1, to=2-2]
	\arrow[hook, from=2-2, to=1-2]
\end{tikzcd}\]
commute with $F$ being continuous. We have to show that $f$ is continuous, but by definiton of the exponential topology this is equivalent to 
\[
Z \times X \to B, \; (z,x) \mapsto f(z)(x)
\]
being continuous and continuity of $F$ is equivalent to continuity of 
\[
Z \times X \to Y, \; (z,x) \mapsto F(z)(x),
\]
so (i) follows. For (ii) note that the restriction map is continuous if and only if 
\[
C^0(X,Y) \times A \to Y, \; (f,x) \mapsto f|_A(x)
\]
is continuous. But this is just the restriction of the evaluation map which is continuous by \cite{steimle_topo}. Similiarly the map in (iii) is continuous if and only if 
\[
C^0(X,Y) \times C^0(X,Z) \times X \to Y \times Z, \; (f,g,x) \mapsto (f(x),g(x))
\]
is. But again the continuity follows from the continuity of the evaluation map.
\end{proof}

\begin{proof}[Proof of Proposition \ref{whitney_top:prop:ck_conv_on_rn}]
    In \cite{golubitsky2012stable} it was shown that for compact domains $C^k$-convergence is equivalent to the $k$-jets converging uniformly with respect to any metric on the Jet-bundle compatible with its topology, that is 
    \[
    f_n \xrightarrow{C^k} f \; \Leftrightarrow \; j^kf_n \xrightarrow{\mathrm{uniformly}} j^kf
    \]
    Note that for compact domains uniform convergence and $C^0_\mathrm{exp}$-convergence agree (see \cite{steimle_topo}). Now let $d(n,k) = \frac{(n+k)!}{n!k!}$ be the total number of all partial derivatives $\partial^\alpha$ of order $1 \leq \abs{\alpha} \leq k$ and let 
    \[
    \tau \colon D \times V \times \R^{d(n,k)} \xrightarrow{\cong} J^k(D,V)
    \]
    be the canonical diffeomorphism. Then 
    \[
    j^kf_n \xrightarrow{C^0_\mathrm{exp}} j^kf \; \Leftrightarrow \; j^kf_n \circ \tau \xrightarrow{C^0_\mathrm{exp}} j^kf \circ \tau.
    \]
    But a sequence of maps with values in a product $C^0_\mathrm{exp}$-converges if and only if each component $C^0_\mathrm{exp}$-converges. Since the components of $j^kf_n \circ \tau$ and $j^kf \circ \tau$ are precisely the partial derivatives of $f_n$ and $f$ up to order $k$ the statement follows (after again applying the equivalence of uniform and $C^0_\mathrm{exp}$-convergence on compact domains).
\end{proof}

\begin{proof}[Proof of Proposition \ref{whitney_top:prop:ck_conv_equivalent_characterizations}]
    Let $K \subseteq X$ be any compact codimension $0$ submanifold. It then has been shown in \cite{golubitsky2012stable} that 
    \[
    f_n|_K \xrightarrow{C^k} f|_K
    \]
    is equivalent to 
    \[
    j^kf_n|_K \xrightarrow{\mathrm{uniformly}}j^kf|_K
    \]
    with respect to any metric on $J^k(X,Y)$ compatible with the topology. But again uniform convergence on compact domains is equivalent to $C^0_\mathrm{exp}$-convergence, which in turn is equivalent to $C^0$-convergence by Lemma \ref{whitney_top:lem:c0_and_exponential_topology_agree}. This shows the equivalence of (i) and (ii). 
    
    For (ii) $\Rightarrow$ (iii) let $\phi,\psi$ be as in (iii) and let $D\subseteq U$ be any compact codimension $0$ submanifold. Then by Lemma \ref{whitney_top:lem:c0_and_exponential_topology_agree} we again observe that $C^0$-convergence of $j^kf_n$ on $\phi(D)$ is equivalent to $C^0_\mathrm{exp}$-convergence and since the exponential topology behaves nicely with respect to restiction in the codomain (see Lemma \ref{whitney_top:lem:exponential_topology_facts}) we get that the restricted maps 
    \[
    j^kf_n|_{\phi(D)} \colon \phi(D) \to J^k(U^\prime,V^\prime)
    \]
    $C^0_\mathrm{exp}$-converge to 
    \[
    j^kf|_{\phi(D)} \colon \phi(D) \to J^k(U^\prime,V^\prime),
    \]
    where $J^k(U^\prime,V^\prime) \subseteq J^k(X,Y)$ is considered an open subset. Under the diffeomorphisms
    \begin{align*}
        D &\xrightarrow{\cong} \phi(D), \\
        U \times V \times \R^{d(n,k)} &\xrightarrow{\cong} J^k(U^\prime,V^\prime)
    \end{align*}
    induced by $\phi$ and $\psi$ the maps $j^kf_n|_{\phi(D)}, j^kf|_{\phi(D)}$ correspond to maps 
    \[
    g_n,g \colon D \to U \times V \times \R^{d(n,k)}
    \]
    and $g_n \xrightarrow{C^0_\mathrm{exp}} g$. But a sequence of maps with values in a product $C^0_\mathrm{exp}$-converges if and only if each component converges (again see Lemma \ref{whitney_top:lem:exponential_topology_facts}). Since the components of $g_n$ are $\mathrm{id}_U$ and the partial derivatives of $\psi^{-1} \circ f_n \circ \phi$ of order less than $k$ (and similiarly for $g$) we conclude that 
    \[
    \frac{\partial^\alpha(\psi^{-1} \circ f_n \circ \phi)}{\partial x^\alpha}\scalebox{1}[1.8]{$|$}_D \xrightarrow{C^0_\mathrm{exp}} \frac{\partial^\alpha(\psi^{-1} \circ f \circ \phi)}{\partial x^\alpha}\scalebox{1}[1.8]{$|$}_D
    \]    
    for all $\alpha \in \N_0^n$ with $0 \leq |\alpha| \leq k$. Now (iii) follows from Proposition \ref{whitney_top:prop:ck_conv_on_rn} (and the equivalence of uniform and $C^0_\mathrm{exp}$-convergence on compact domain). 

    Clearly (iii) $\Rightarrow$ (iv), so let's assume (iv). Using the same notation and argument as above but reversed we see that 
    \[
    (\psi^{-1} \circ f_n \circ \phi)|_D \xrightarrow{C^k} (\psi^{-1} \circ f \circ \phi)|_D
    \]
    if and only if 
    \[
    j^kf_n|_{\phi(D)} \xrightarrow{C^0_\mathrm{exp}} j^kf|_{\phi(D)},
    \]
    where we consider $j^kf_n|_{\phi(D)}$ and $j^kf|_{\phi(D)}$ as maps 
    \[
    \phi(D) \to J^k(X,Y).
    \]
    Now choose any metric on $J^k(X,Y)$ compatible with the topology. Then $C^0_\mathrm{exp}$-convergence of $j^kf_n|_{\phi(D)}$ is equivalent to uniform convergence with respect to that metric. Thus we have shown that around each point $p \in X$ theres exists some neighbourhood $\phi(D)$ of $p$, such that $j^kf_n$ converges uniformly to $j^kf$ on $\phi(D)$. But any compact codimension $0$ submanifold $K\subseteq X$ is covered by finitely many such neighbourhoods and it is easily seen that this already implies uniform convergence of $j^kf_n$ to $j^kf$ on all of $K$. Then again using the equivalence of uniform convergence and $C^0_\mathrm{exp}$-convergence on compact domains as well as Lemma \ref{whitney_top:lem:c0_and_exponential_topology_agree} we conclude that 
    \[
    j^kf_n \xrightarrow{C^0_\mathrm{loc}} j^kf.
    \]
\end{proof}